\subsection{Clifford'o algebros}


Toliau nagrinėsime nagrinėsime dvi Clifford'o algebras $\Clp$ ir $\Clm$.
Tai yra vektorinės erdvės su papildoma bitiesine daugybos operacija,
kuri leidžia sudauginti du vektorius ir gauti naują vektorių.
Šias algebras apibrėšime pirma paimdami daug bendresnę tenzorinę algebrą
ir tada jai suteiksime norimą struktūrą.

\begin{definition}
	Vienetinė tenzorinė algebra $T(\set R^3)$ yra tiesioginė suma
	\begin{equation}
		T(\mathbb{R}^3) 
		=
		\bigoplus_{k=0}^{\infty} T^k
		=
		\mathbb{R} 
		\oplus \set R^3 
		\oplus (\set R^3 \ot \set R^3) 
		\oplus (\set R^3 \ot \set R^3 \ot \set R^3) 
		\oplus \dots
		\,,
	\end{equation}
	$T^k$ čia yra vektorinės erdvės $T^0 = \set R$\: ,\: $T^1=\set R^3$\: ,\: $T^2=\set R^3 \ot \set R^3$\: ,\:
	$T^3=\set R^3 \ot \set R^3 \ot \set R^3$
	ir taip toliau.
	Algebros daugyba yra duodama vektorių tenzorinės sandaugos.
	Tegul $\mathbb1 \in T^0$ žymi vienetinį tenzorinės algebros elementą,
	kuris su visais $v \in T(\set R^3)$ tenkina
	$\mathbb1 \otimes v = v \otimes \mathbb1 = v$.
\end{definition}
Kitaip pasakius tai yra algebra,
kurios vektoriai yra įvairių ilgiu tenzorinės sandaugos
$\vb v_1 \ot \vb v_2 \ot \dots \ot \vb v_n$\:,\: $\vb v_i \in \set R^3$
ir tokių sandaugų tiesinės kombinacijos.
\begin{example}
	Vektorius $3\;\mathbb 1 + \vb v$ priklauso $T(\set R^3)$
	ir yra suma vektorių iš $T^0$ ir $T^1$.
	Vektorius $2\vb w + \vb v \ot \vb w$ yra suma vektorių iš $T^1$ ir $T^2$.
	Jų sandauga tenzorinėje algebroje yra
	\begin{equation}
		\begin{split}
			\qty(3\;\mathbb 1 + \vb v)
			\ot
			\qty(2\vb w + \vb v \ot \vb w)
			=&
			3\;\mathbb 1 \ot 2\vb w
			+ 3\;\mathbb 1 \ot \vb v \ot \vb w
			+ \vb v \ot 2\vb w
			+ \vb v \ot \vb v \ot \vb w
			\\=&
			3\cdot2\qty(\mathbb 1 \ot \vb w)
			+ 3\qty(\mathbb 1 \ot \vb v \ot \vb w)
			+ 2(\vb v \ot \vb w)
			+ \qty(\vb v \ot \vb v \ot \vb w)
			\\=&
			6 \vb w
			+ 3\qty(\vb v \ot \vb w)
			+ 2(\vb v \ot \vb w)
			+ \qty(\vb v \ot \vb v \ot \vb w)
			\\=&
			6 \vb w
			+ 5\qty(\vb v \ot \vb w)
			+ \qty(\vb v \ot \vb v \ot \vb w)
			\,,
		\end{split}
	\end{equation}
	čia pirmoje eilutėje pasinaudojome tuo kad tenzorinės sandauga
	distributyvi; antroje eilutėje pasinaudojame tuo kad tenzorinė sandauga
	tiesiška kiekvieno dauginamo vektoriaus atžvilgiu, kad iškeltume 
	skaliarinius daugiklius; trečioje eilutėje sudedame tiesiškai
	priklausomus vektorius.
	Gautas vektorius yra suma vektorių iš $T^1$, $T^2$ ir $T^3$.
\end{example}


Dabar $T(\set R^3)$ suteiksime papildomą struktūrą, kad algebroje
$\Clp$ būtų tenkinamas sąryšis
\begin{equation}
	\vb v \ot \vb w
	+
	\vb w \ot \vb v
	=
	2(\vb v , \vb w) \mathbb 1
	\quad ,
	\label{eq:cl+}
\end{equation}
o algebroje $\Clm$ tenkinamas sąryšis
\begin{equation}
	\vb v \ot \vb w
	+
	\vb w \ot \vb v
	=
	- 2(\vb v , \vb w) \mathbb 1
	\quad .
	\label{eq:cl-}
\end{equation}
Šiuos sąryšius ,,užkoduojame'' į aibes $S^\pm \subset T(\set R^3)$
\begin{equation}
	\begin{split}
		S^+ = \qty{
			\vb v \ot \vb w
			+
			\vb w \ot \vb v
			-
			2(\vb v , \vb w) \mathbb 1
			\mid
			\vb v, \vb w \in \set R^3
		}
		\\
		S^- = \qty{
			\vb v \ot \vb w
			+
			\vb w \ot \vb v
			+
			2(\vb v , \vb w) \mathbb 1
			\mid
			\vb v, \vb w \in \set R^3
		}\,.
	\end{split}
\end{equation}
Užrašome tenzorinės algebros poalgebrius $I^\pm \subset T(\set R^3)$,
gaunamus iš visų vektorių, kurie turi bent vieną daugiklį iš $S^\pm$
\begin{equation}
	I^\pm = \qty{
		\sum_{\substack{
			s \in S^\pm \\ 
			a,b \in \mathup{T(V)} 
		}}
		a \ot s \ot b
	}
	\quad .
\end{equation}
Šie poalgebriai papildomai tenkina savybę, kad bet kokiems $a \in T(\set R^3)$
ir $d \in I^\pm$, $a\ot d, d \ot a \in I^\pm$. Poalgebris
tenkinantis šia savybę yra vadinamas algebros idealu.
Literatūroje šią konstrukciją yra įprasta pateikti sakiniu:
,,$I^\pm$ yra algebros $T(\set R^3)$ idealas generuojamas 
sąryšio (\ref{eq:cl+}) ar atitinkamai (\ref{eq:cl-})''.

Duotam vektoriui $a \in T(\set R^3)$ aibė
$\qty{a+d \mid d \in I^\pm}$ (dar žymima $a+I^\pm$)
vadinama gretutine klase (ang. \emph{coset} arba \emph{equivalence class}),
o vektorius $a$ vadinamas klasės atstovu.
Dviejų gretutinių klasių $a+I^\pm$ ir $b + I^\pm$ sandauga yra gretutinė
klasė, kuriai priklauso atstovų sandauga $a \ot b$.
Ši sandauga nepriklauso nuo pasirinktų klasės atstovų, nes
\begin{equation}
	(a+I^\pm)\ot(b+I^\pm)
	=
	a \ot b
	+ a \ot I^\pm
	+ I^\pm \ot b
	+ I^\pm \ot I^\pm
	=
	a \ot b + I^\pm
	\,.
\end{equation}
Sumos nariai 
$a \ot I^\pm = I^\pm $,
$I^\pm \ot b = I^\pm $,
$I^\pm \ot I^\pm = I^\pm $,
nes $I^\pm$ yra algebros $T(\set R^3)$ idealas.
Tai tikrai
\begin{equation}
	(a+I^\pm)\ot(b+I^\pm)
	=
	a \ot b + I^\pm
	\,.
\end{equation}



\begin{definition}
	Algebra $\Clpm$ yra faktorinė algebra $T(\set R^3)/I^\pm$,
	kurios elementai yra gretutinės klasės $a+I^\pm$, $a \in T(\set R^3)$
	ir jų sandauga yra $(a+I^\pm)\ot(b+I^\pm)=a\ot b + I^\pm.$
\end{definition}

\begin{example}
	$\vb w \ot \vb v +\vb v \ot \vb w + I^+$ yra gretutinė klasė,
	kuri yra vektorius algebroje $\Clp$. 
	Supaprastiname išraišką
	\begin{equation}
		\vb w \ot \vb v +\vb v \ot \vb w + I^+
		= 
		\vb w \ot \vb v +\vb v \ot \vb w 
		- \qty(
			\vb v \ot \vb w + \vb w \ot \vb v 
			- 2(\vb v,\vb w)\mathbb 1
		) 
		+ I^+
		=
		2(\vb v,\vb w)\mathbb 1 + I^+
		\,,
	\end{equation}
	čia pasinaudojome tuo kad 
	$\qty( \vb v \ot \vb w + \vb w \ot \vb v - 2(\vb v,\vb w)\mathbb 1) \in I^+$.
	Pavyko parodyti, kad $\vb w \ot \vb v + \vb v \ot \vb w$ 
	ir $2(\vb v,\vb w) \mathbb 1$ atstovauja tą pačią klasę.
\end{example}

\begin{example}
	Tegul $a,b \in T(\set R^3)$ ir
	$ a \ot (\vb v \ot \vb w + 2(\vb v,\vb w)\mathbb 1) \ot b + I^- $
	yra gretutinė klasė, kuri yra vektorius iš $\Clm$. Šią išraišką supaprastiname atlikdami veiksmus
	\begin{equation}
		\begin{aligned}
			a \ot (\vb v \ot \vb w + 2(\vb v,\vb w)\mathbb 1) \ot b + I^- 
			& = a \ot (\vb v \ot \vb w + 2(\vb v,\vb w)\mathbb 1) \ot b
			\\& \quad - a \ot \qty(\vb v \ot \vb w + \vb w \ot \vb v +2(\vb v,\vb w) \mathbb 1)  \ot b + I^-
			\\& = a \ot \qty( \vb v \ot \vb w + 2(\vb v,\vb w)\mathbb 1 - \vb v \ot \vb w - \vb w \ot \vb v -2(\vb v,\vb w) \mathbb 1))\ot b + I^-
			\\& = a \ot \qty( - \vb w \ot \vb v )\ot b + I^-
			\\& = - a \ot	\vb w \ot \vb v	\ot b + I^- 
			\,.
		\end{aligned}
	\end{equation}
	Parodėme, kad sudarius faktorinę algebrą $\Clm$,
	galime jos viduje taikyti sąryšį (\ref{eq:cl-})
	manipuliuojant gretutinės klasės atstovus.
\end{example}
Toliau priimame sutrumpintą žymėjimą, nes išreikštai
su gretutinėmis klasėmis nebedirbsime. Užrašysime tik
klasės atstovą, o kad išraiška skirtųsi nuo 
vektorių tenzorinės sandaugos vietoje ženklo ,,$\ot$''
tiesiog rašysime vektorius vieną šalia kito.
Vietoje
\begin{equation}
	a \ot (\vb v \ot \vb w + \vb v \ot \vb w) \ot b + I^-
	=
	-2(\vb v,\vb w) \qty(a \ot b) + I^-
\end{equation}
rašysime
\begin{equation}
	a (\vb v \vb w + \vb v \vb w) b
	=
	-2(\vb v,\vb w) a b
	\,.
\end{equation}

\begin{proposition}
	\label{prop:R_action}
	Galioja tapatybės:
	\begin{enumerate}
		\item[$(i)$] algebrose $\Clpm$ turime $ \vb v \vb v = \pm(\vb v,\vb v) \mathbb 1$
		\item[$(ii)$] algebroje $\Clp$ turime $ -\vb v \vb w \vb v = -2(\vb v,\vb w) \vb v + (\vb v,\vb v)\vb w $
		\item[$(iii)$] algebroje $\Clm$ turime $ \vb v \vb w \vb v = -2(\vb v,\vb w) \vb v + (\vb v,\vb v)\vb w $
	\end{enumerate}
\end{proposition}
\begin{proof} Taikydami sąryšį $\vb v \vb w + \vb w \vb v = \pm2(\vb v ,\vb w)$ gauname
	\begin{enumerate}
		\item[$(i)$] 
			$
			\vb v \vb v 
			= \frac{1}{2}(\vb v \vb v + \vb v \vb v) 
			= \pm(\vb v,\vb v) \mathbb 1
			\,,
			$
		\item[$(ii)$] 	
			$
			-\vb v \vb w \vb v 
			= -\qty( 2(\vb v,\vb w)\mathbb 1 - \vb{wv} )\vb v 
			= -2(\vb v,\vb w) \vb v + \vb{wvv} 
			= -2(\vb v,\vb w) \vb v + (\vb v,\vb v)\vb w 
			\,,
			$
		\item[$(iii)$]
			$
			\vb v \vb w \vb v 
			= \qty( - 2(\vb v,\vb w)\mathbb 1 - \vb{wv} )\vb v 
			= -2(\vb v,\vb w) \vb v - \vb{wvv} 
			= -2(\vb v,\vb w) \vb v + (\vb v,\vb v)\vb w 
			\,.
			$
			\qedhere
	\end{enumerate}
\end{proof}

\begin{remark}
	Atkreipiame dėmesį, kad kai $(\vb v,\vb v)=1$,
	$(\vb v,\vb w)\vb v$ yra vektoriaus $\vb w$ projekcija į $\vb v$
	ir $-2(\vb v,\vb w)\vb v + \vb w$ yra vektoriaus atspindys
	per plokštumą statmeną vektoriui $\vb v$. 
\end{remark}

Tegul $\qty{\vb e_1,\vb e_2,\vb e_3}$ yra ortonormali $\set R^3$ bazė.
Bazinius vektorius galime perkelti į algebra $\mathup{Cl^\pm(V)}$,
kur jie turi tenkinti
\begin{equation}
	\begin{split}
		\vb e_i \vb e_j &= - \vb e_j \vb e_i \ , \ i \neq j \\
		\vb e_i \vb e_i &= \pm \mathbb 1
		\quad .
	\end{split}
	\label{eq:base-relations}
\end{equation}
Su šių sąryšių pagalba galime bet kokį $\mathup{Cl^\pm(V)}$ elementą užrašyti,
kaip tiesinę kombinaciją aštuonių elementų:
\begin{equation}
	\mathbb 1 \,,\quad
	\vb e_1 \,,\quad 
	\vb e_2 \,,\quad  
	\vb e_3 \,,\quad  
	\vb e_2 \vb e_3 \,,\quad
	\vb e_1 \vb e_3 \,,\quad
	\vb e_1 \vb e_2 \,,\quad 
	\vb e_1 \vb e_2 \vb e_3 \,,
\end{equation}
nes su pirma taisykle bet kurį tenzorinės algebros bazės elementą galima išrikiuoti,
taip kad ortų indeksai būtų surašyti didėjimo tvarka
ir su antra taisykle pasikartojantys ortai panaikinami.
Tai reiškia $\Clpm$ yra aštuonmetės algebros.
Įvedame pažymėjimą, kurį toliau naudosime algebrų $\Clpm$ bazei
\begin{equation}
	\vb e_{23} := \vb e_2 \vb e_3
	\,,\quad
	\vb e_{31} := \vb e_3 \vb e_1 = -\vb e_1 \vb e_3 
	\,,\quad
	\vb e_{12} := \vb e_1 \vb e_2
	\,,\quad
	\vb e_{123} := \vb e_1 \vb e_2 \vb e_3
	\,.
\end{equation}

\begin{definition}
	Transponuota $\Clpm$ bazė yra 
	\begin{equation}
		\qty(\vb e_{i_1} \vb e_{i_2} \dots \vb e_{i_k} )^t
		:=
		\vb e_{i_k} \vb e_{i_{k-1}} \dots \vb e_{i_1}
		=
		(-1)^\frac{k(k-1)}{2} \vb e_{i_1} \vb e_{i_2} \dots \vb e_{i_k} 
	\end{equation}
	ir transponavimo operacija yra tiesiškai praplečiama
	bet kuriam algebros $\Clpm$ vektoriui, veikimu ant bazės.
	Transponavimas yra algebros anti-automorfizmas,
	tai reiškia, kad bet kuriems algebros elementams $a$ ir $b$ galioja
	\begin{equation}
		(a b)^t = b^t a^t
		\,.
	\end{equation}
\end{definition}
\begin{definition}
	Apibrėžiame algebrų $\Clpm$ automorfizmą $\alpha$, jo veikimu ant algebrų bazės
	\begin{equation}
		\alpha( \vb e_{i_1} \vb e_{i_2} \dots \vb e_{i_k} )
		:=
		(-1)^{k} \vb e_{i_1} \vb e_{i_2} \dots \vb e_{i_k} 
		\,.
	\end{equation}
\end{definition}
Nesunku pamatyti,
kad nėra svarbu kokia tvarka pritaikomas transponavimas ir automorfizmas $\alpha$,
tai $\alpha(v)^t = \alpha(v^t)$.


Pažymime algebros poerdvius
\begin{equation}
	\begin{split}
		A_0^\pm &:= \qty{
			v \in \Clpm \mid \alpha(v)=v
		}
		\,,
		\\
		A_1^\pm &:= \qty{
			v \in \Clpm \mid \alpha(v)=-v
		}
		\,.
	\end{split}
\end{equation}
Susiejant su erdvės baze šiuos poerdvius galima užrašyti kaip
tiesinį apvalką pasirinktų bazės elementų
\begin{equation}
	\begin{aligned}
		A_0^\pm &= \mathup{span}\qty(
			\mathbb 1,\,
			\vb e_{23},\,
			\vb e_{31},\,
			\vb e_{12}
		)
		\,,
		\\
		A_1^\pm &= \mathup{span}\qty(
			\vb e_1 ,\,
			\vb e_2 ,\,
			\vb e_3 ,\,
			\vb e_{123}
		)
		\,.
	\end{aligned}
\end{equation}

\begin{example}
	Tegul $v_0, v'_0 \in A_0^\pm$ ir $v_1, v'_1 \in A_1^\pm$. Turime
	\begin{equation}
		\begin{split}
			\alpha(v_0 v'_0)
			= \alpha(v_0)\alpha(v'_0)
			= v_0 v'_0
			\implies &
			v_0 v'_0 \in A_0^\pm
			\,,
			\\
			\alpha(v_0 v_1)
			= \alpha(v_0)\alpha(v_1)
			= v_0 (-v_1)
			= - v_0 v_1
			\implies &
			v_0 v_1 \in A_1^\pm
			\,,
			\\
			\alpha(v_1 v_0)
			= \alpha(v_1)\alpha(v_0)
			= (- v_1)v_0
			= - v_1 v_0
			\implies &
			v_1 v_0 \in A_1^\pm
			\,,
			\\
			\alpha(v_1 v'_1)
			= \alpha(v_1)\alpha(v'_1)
			= (- v_1)(- v'_1)
			= v_1 v'_1
			\implies &
			v_1 v'_1 \in A_0^\pm
			\,.
		\end{split}
	\end{equation}
	Taip yra tokia pati struktūra kaip grupė 
	$\set Z_2 \cong C_2$ su sudėtimi
	\begin{equation}
		\begin{split}
			\mathbf{lyginis} + \mathbf{lyginis} &= \mathbf{lyginis}
			\,,
			\\
			\mathbf{lyginis} + \mathbf{nelyginis} &= \mathbf{nelyginis}
			\,,
			\\
			\mathbf{nelyginis} + \mathbf{lyginis} &= \mathbf{nelyginis}
			\,,
			\\
			\mathbf{nelyginis} + \mathbf{nelyginis} &= \mathbf{lyginis}
			\,.
		\end{split}
	\end{equation}
\end{example}
Dėl tokios daugybos struktūros vektoriai erdvėje $A_0$ yra vadinami \emph{lyginiais},
o vektoriai erdvėje $A_1$ yra vadinami \emph{nelyginiais}.
Pamatome, kad automorfizmas $\alpha$ 
leidžia algebrą $\Clpm$ suskaidyti 
į lyginio ir nelyginio poerdvių tiesioginę sumą $A_0 \oplus A_1$.
Šis algebros išskaidymas į vektorinius poerdvius vadinamas
,,$\set Z_2$ gradavimu'', kurį suteikia automorfizmas $\alpha$.
Lyginis poerdvis yra uždaras algebros daugybos atžvilgiu,
todėl tai yra $\Clpm$ poalgebris.

\begin{proposition}
	\label{prop:even_iso}
	Algebrų $\Clp$ ir $\Clm$
	Lyginiai poalgebriai $A^+_0$ ir $A^-_0$ yra izomorfiškos algebros.
\end{proposition}

\begin{proof}
Algebros $A_0^+$ ir $A_0^-$ turi bazę $\qty{
		\mathbb 1,\, 
		\vb e_{23},\,
		\vb e_{31},\,
		\vb e_{12}
	}$. 
Pastebime kad algebroje $A_0^+$
\begin{gather}
	\vb e_{23} \vb e_{31} 
	= \vb e_2 \vb e_3 \vb e_3 \vb e_1
	= \vb e_2 \vb e_1
	= - \vb e_{12}
	\,,
\\
	\vb e_{31} \vb e_{23} 
	= \vb e_3 \vb e_1 \vb e_2 \vb e_3
	= \vb e_1 \vb e_3 \vb e_3 \vb e_2
	= \vb e_1 \vb e_2
	= \vb e_{12}
	\,,
\\
	{\vb e_{31}}^2 
	= \vb e_3 \vb e_1 \vb e_3 \vb e_1 
	= - \vb e_3 \vb e_3 \vb e_1 \vb e_1 
	= - \mathbb 1
	\,,
\\
	{\vb e_{23}}^2 
	= \vb e_2 \vb e_3 \vb e_2 \vb e_3 
	= - \vb e_2 \vb e_2 \vb e_3 \vb e_3 
	= - \mathbb 1
	\,,
\end{gather}
o algebroje $A_0^-$
\begin{gather}
	\vb e_{23} \vb e_{31} 
	= \vb e_2 \vb e_3 \vb e_3 \vb e_1 
	= - \vb e_2 \vb e_1 
	= \vb e_1 \vb e_2 
	= \vb e_{12}
	\,,
\\
	\vb e_{31} \vb e_{23} 
	= \vb e_3 \vb e_1 \vb e_2 \vb e_3
	= \vb e_1 \vb e_3 \vb e_3 \vb e_2
	= - \vb e_1 \vb e_2
	= - \vb e_{12}
	\,,
\\
	{\vb e_{31}}^2 
	= \vb e_3 \vb e_1 \vb e_3 \vb e_1 
	= - \vb e_3 \vb e_3 \vb e_1 \vb e_1 
	= - \mathbb 1
	\,,
\\
	{\vb e_{23}}^2 
	= \vb e_2 \vb e_3 \vb e_2 \vb e_3 
	= - \vb e_2 \vb e_2 \vb e_3 \vb e_3 
	= - \mathbb 1
	\,.
\end{gather}
Abiems algebroms $\vb e_{31}$ ir $\vb e_{23}$ yra algebros generatoriai.
Jie tenkina sąryšius
\begin{equation}
	\begin{aligned}
		\vb e_{31} \vb e_{23} &= - \vb e_{23} \vb e_{31} \,,\\
		{\vb e_{31}}^2 &= - \mathbb1 \,,\\
		{\vb e_{23}}^2 &= - \mathbb1 \,.
	\end{aligned}
\end{equation}

Tegul $f:A^+_0 \to A^-_0$ yra tiesinis atvaizdavimas,
apibrėžiamas veikimu ant $A^+_0$ bazės
\begin{equation}
	\mathbb1 \mapsto \mathbb1
	\,,\quad
	\vb e_{23} \mapsto -\vb e_{23}
	\,,\quad
	\vb e_{31} \mapsto -\vb e_{31}
	\,,\quad
	\vb e_{12} \mapsto -\vb e_{12}
	\,.
\end{equation}
Akivaizdu kad $f$ yra bijekcija.
Įrodymui kad $f$ yra algebrų homomorfizmas užtenka parodyti,
kad $f$ nepakeičia generatorių daugybos sąryšių
\begin{equation}
	\begin{aligned}
		f(\vb e_{31}) f(\vb e_{23}) 
		=(-\vb e_{31})(-\vb e_{23}) 
		&= - \vb e_{23} \vb e_{31} 
		= - f(\vb e_{23}) f(\vb e_{31}) 
\,,\\
		f(\vb e_{31})^2 
		=(-\vb e_{31})^2 
		&= - \mathbb1
		= - f(\mathbb1)
\,,\\
		f(\vb e_{23})^2 
		=(-\vb e_{23})^2 
		&= - \mathbb1
		= - f(\mathbb1)
\,.
	\end{aligned}
\end{equation}
Atvaizdavimas $f$ yra algebrų izomorfizmas,
tad $A^+_0 \cong A^-_0$.
\end{proof}
